% !TEX root = ../main.tex
\section{Discussion}
\label{sec:discussion}

This paper makes four connected points. First, the Hamiltonian of mean force is
best viewed as a constructive reduced generator, not merely as a formal
logarithm. Second, in the Gaussian sector this constructive question reduces to
a closure problem for the adjoint chain generated by \(H_Q\) and the coupling
operator. Third, the same structure can be written as an explicit
compositional architecture, which is why the KAN comparison can be made
formally rather than rhetorically. Fourth, the same closure logic explains both
exact solvability and the failure modes of finite-basis numerics.

\subsection*{Closure and constructive reduced generators}

The central algebraic result is that representability of the HMF inside a
restricted operator family has an exact answer. Conditions (C1)--(C3) separate
the problem into three layers: whether repeated adjoint action grows the
operator basis, whether the Gaussian influence remains inside that basis, and
whether the BCH recombination generates further operators. When all three hold,
the reduced equilibrium admits an exact finite generator. When they fail, the
failure is structural rather than computational: the reduced state can still be
computed, but not compressed losslessly into the proposed ansatz.

This makes the closure theorem useful beyond the present example. It turns the
usual question ``can we evaluate the reduced state?'' into the sharper question
``which operator family is large enough to contain the exact reduced
generator?'' In that form, the theorem is already a guide to controlled
approximation. One may truncate the operator side without approximating the bath
side, and the resulting error has a clear interpretation as model-class error.

\subsection*{The qubit as an exact noncommuting benchmark}

For the spin-boson qubit, closure is exact. The reduced generator remains in
the Bloch plane determined by the system axis and the coupling axis, and the
entire bath compresses to two scalar response channels together with the
influence magnitude \(\chi\). To our knowledge, this makes
Eq.~\eqref{eq:HMF_qubit_closed_hmf_ai} the first closed-form Hamiltonian of
mean force for a genuinely noncommuting open quantum system. Earlier exact
benchmarks are either commuting, quadratic/Gaussian, or asymptotic. Here none
of those simplifications is doing the work. The work is done by complete
closure of the compositional architecture itself.

That mechanism is worth isolating. The adjoint orbit terminates into a
two-family alternation, the bilinear influence remains in the Pauli algebra,
and the nonlinear BCH tower resums because the traceless influence obeys
$M^2=\chi^2\mathbb I$. Figure~\ref{fig:bloch_overview} is the geometric
signature of that fact: bare state, influence state, and final reduced state
all remain inside one Bloch plane. What should generalize to larger systems is
not the specific Pauli formulas but the separation between bath data and
operator growth, together with the question of whether every layer of that
structure still closes.

\subsection*{Representability bottlenecks and finite-resource inference}

The numerical diagnostics sharpen the interpretation. ED disagrees with the
analytic prediction not when the theory becomes inaccurate, but when the chosen
finite basis is no longer expressive enough to realize the exact dressed state.
The response scale \(\chi\sim 1\) marks the onset of this bottleneck in the
examples studied here. In other words, the crossover is simultaneously physical
and representational: it tracks both the reorganization of the reduced state
and the point at which a restricted model class begins to fail.

This is why the distinction between discrete analytic and numerical ED matters.
Once the discrete analytic branch peels away from ED and returns toward the
continuum benchmark, the problem is no longer one of physics missed by the
generator. It is one of structure missed by the basis. That diagnosis is useful
well beyond this example, because many approximate methods in open-system
physics are precisely choices of tractable operator or Hilbert-space classes.

\subsection*{Compositional architectures, representation cost, and KANs}

Recent work on representation cost emphasizes that exact state descriptions are
less useful than explicit generators that expose what must be retained under
coarse-graining~\cite{mccaulCostComplexity2026}. The HMF offers a concrete
equilibrium realization of that principle. The partial trace hides a large
environment inside a reduced state; closure determines whether that hidden
structure can still be recast as a compact generator.

The strengthened KAN correspondence is not decorative. Kolmogorov--Arnold
networks ask when multivariate dependence can be realized by a compact
composition of univariate edge functions and additive aggregation
~\cite{liuKANKolmogorovArnoldNetworks2024}. The Gaussian HMF problem already
has that architecture built in. The adjoint orbit is the univariate basis in
operator space, the ordered-kernel matrix $C^>_{nm}(\beta)$ is the aggregation
tensor, and the Todd/Bernoulli BCH resummation is the nonlinear composition
rule. Here those objects are not learned by optimization; they are supplied in
closed form by equilibrium thermodynamics.

In that dictionary, the operator ansatz $\mathcal A$ is the model class and
closure failure is model-class insufficiency. The exact reduced generator still
exists, but a fixed finite representation cannot realize it without projection.
That is why the non-quadratic continuous-variable obstruction and the
finite-basis numerical bottleneck belong in the same discussion: both are
instances of a universal target outgrowing the chosen compact representation.
Mechanistic interpretability enters only after that point. Once the layers are
explicit, one can inspect which channels carry the reduced generator and which
ones force basis growth.

\subsection*{Limitations and outlook}

The main analytic restriction is Gaussianity. The quenched representation itself
is more general, but the exact collapse of the influence to a bilinear operator
depends on the bath cumulant hierarchy truncating at second order. Beyond that
sector, higher cumulants generate higher operator monomials and the closure
problem becomes correspondingly richer.

Within the Gaussian sector, the next natural targets are multi-qubit systems and
other finite-dimensional models where the adjoint chain generates nontrivial
many-body operator content. There the central question is no longer whether an
HMF exists, but how rapidly operator growth escapes a chosen locality class. In
that sense the present qubit solution is a starting point for a broader study
of constructive reduced generators in interacting systems. A second direction
is methodological: once the reduced generator is written as an explicit
compositional object, one can ask which approximate schemes preserve that
architecture and which ones merely fit the reduced state after the fact.
